% Archivo: 1_intro.tex
\section{Introducción}

El estudio del movimiento de un gran número de partículas es fundamental en la física estadística y computacional. En este trabajo, se analizan tres fenómenos interconectados: los caminos aleatorios (random walks) en múltiples dimensiones, la difusión macroscópica y microscópica de un soluto (modelo de café y leche), y la percolación de invasión en medios porosos. 

El objetivo principal de esta práctica es demostrar cómo las fluctuaciones estocásticas microscópicas dan lugar a comportamientos macroscópicos deterministas. Se busca validar numéricamente la relación lineal entre el desplazamiento cuadrático medio y el tiempo en caminos aleatorios, comprobar la equivalencia entre la resolución de la ecuación de difusión mediante diferencias finitas y la evolución entrópica de un enjambre de partículas, y finalmente, simular el avance de un fluido invasor a través de un medio con resistencia aleatoria, minimizando la energía capilar en cada paso. El alcance computacional incluye simulaciones en 1D, 2D y 3D, el uso de esquemas numéricos explícitos y la implementación de colas de prioridad para optimizar la percolación.